\section{Threats to validity}
\label{sec:threats-to-validity}
%Threats to Validity discussed (Petersen, Vakkalanka, \& Kuzniarz, 2015; Wohlin et al., 2012); Outlines key benefits, challenges, and research opportunities.

%\DD{}{Structure the section identifying the categories of the main treats that apply to our study}

%\DD{}{For each threat, try to highlight the way we tried to mitigate or address it, if any.}

\begin{comment}
The validity of our systematic review is subject to several interrelated threats that range different phases of the research process, including literature search, study selection, quality assessment, and data synthesis. Petersen et al. in \cite{PETERSEN20151} note that an inadequately formulated search strategy jeopardizes external validity by potentially omitting relevant literature; indeed, although we broadened our search string with relaxed boundaries and multiple synonymous terms, there remains a risk that pertinent studies were overlooked—further compounded by publication bias (e.g., the file-drawer effect) and language bias, as studies in non-indexed or non-English sources might not have been captured. Similarly, as Wohlin et al. in \cite{wohlin2012} define internal validity as the degree to which causal claims are justified, our parallel selection process, despite leveraging collective expertise and employing deterministic approaches, remains vulnerable to human error, subjective judgment, and comparative bias, where the sequence of review influences inclusion decisions rather than a strictly standardized protocol. Moreover, our quality assessment criteria \DD{}{update according to the final list of QCs}(QC1, QC5, QC7, and QC8) are particularly susceptible to reputation and comparative biases, with high-ranked venues potentially receiving inflated scores and subjective interpretations influencing scoring consistency. In addition, our methodology required the normalization of heterogeneous information across studies to extract unified insights, and these necessary approximations may introduce further inaccuracies, thereby propagating pre-existing biases or mistakes embedded in the original research. Reliability is another concern, as the consistency of data extraction and quality scoring may vary across researchers, even with rigorous inter-researcher discussions, and any divergence in interpretation can affect the reproducibility of our findings. Finally, our reliance on the accuracy of original authors' reporting means that any inherent biases or errors in the primary studies are inevitably carried forward into our synthesis. Despite these limitations, our study has strived to align with the rigorous methodological frameworks advocated by \cite{PETERSEN20151} and \cite{wohlin2012} by establishing explicit definitions, employing standardized scales, and conducting iterative validation sessions to enhance objectivity and reproducibility—though complete elimination of subjectivity remains an ongoing challenge that future research might address through more automated, data-driven evaluation processes.
    
\end{comment}

The validity of our systematic review is subject to several interrelated threats across the different phases of the research process, including literature search, study selection, quality assessment, data extraction, and data synthesis. Following the guidelines of Petersen et al. \cite{PETERSEN20151}, we organized potential threats into five dimensions: descriptive, theoretical, generalizability, interpretive, and repeatability. 
For every dimension, we identify the specific risks that could distort our findings, describe the concrete mitigation actions embedded in our protocol, and acknowledge the residual limitations that remain despite those safeguards.

\textit{Theoretical validity: } Our evidence base may still miss relevant studies or misjudge their merit. File-drawer effects, language bias, and general publication bias can hide pertinent work, while venue reputation may inflate quality scores. Moreover, comparative bias can arise if the sequence in which papers are screened influences inclusion decisions. To mitigate these threats, we broadened the search string with synonyms using the PICOC framework and applied relaxed operators across multiple digital libraries, randomized the order of papers for each reviewer. Screening was performed in parallel, with disagreements resolved by consensus, and quality was assessed through explicit rubrics that we refined until full agreement.

\textit{Descriptive validity: }A significant concern is whether the information we extracted from each primary study accurately reflects what the original authors reported. When compiling the data for our research questions, mapping extracted data to a common field could have introduced errors; for example, when tabulating features, the mapping was subject to researcher interpretation. We mitigated this risk through a shared, version-controlled data extraction form that tracks each mapping to the extracted original source, so that choices could be traced, discussed, and corrected in case of need.

\textit{Generalizability validity: } While this review offers valuable insights into various aspects of customer churn prediction in non-contractual retail settings---including commonly used features, AI techniques, and evaluation metrics---the generalizability of some findings may be limited. The conclusions drawn from the analyzed studies are influenced by specific dataset characteristics, sector-specific dynamics, geographic context, and methodological choices such as feature engineering and model evaluation criteria. Therefore, claims regarding the effectiveness of particular techniques, the relevance of specific features, or the prevalence of certain metrics should be interpreted with caution when applied to different retail domains or operational contexts that are not represented in the reviewed literature.

\textit{Repeatability validity: } For another team to replicate our study, the entire chain of evidence must be visible and reliable. We therefore provide details on the review protocol design and execution (Section \ref{sec:methodology}) and make the data extraction form along with the collected data publicly accessible. 
%(found in Section \ref{sec:methodology} and in \ref{Appendix}), with dates, exact search strings—together with the selection process, the quality-control rubric, and the Data Extraction Form spreadsheet under version control.  

\textit{Residual limitations: } Despite the mitigation techniques adopted above, our search may still overlook some non-indexed studies, and quality scoring retains an element of human subjectivity. By publishing the full protocol and data extraction form, we keep these residual risks transparent and open to scrutiny.
